% !TEX root = ../nc-cubic-non-as.tex
\section{反標準環と二次包絡}
\label{sec:quadratic-obstruction}

\subsection{次数一生成と環論的性質}
\label{subsec:generation-all-degrees}

$\mathcal V=\mathcal O_Z^{\oplus2}\oplus\mathcal O_Z(f)$の大域切断
\[
 x=(1,0,0),\qquad y=(0,1,0),\qquad
 z=(0,0,s),\qquad w=(0,0,t)
\]
を，$\mathcal V\hookrightarrow\mathcal P$を通して$R_1$の元とみなす．
式\eqref{eq:underlying-bundles}より，これらは$R_1$の基底である．

\begin{proposition}
\label{prop:generation-all-degrees}
$R$ は次数一の四元 $x,y,z,w$ で生成される．さらに $R$ は
有限生成可換中心部分環上の有限加群であり，左右 noetherian な domain である．
\end{proposition}

\begin{proof}
中心の切断環を
\[
 S=\bigoplus_{m\ge0}H^0(Z,L^m),\qquad
 \deg H^0(Z,L^m)=2m
\]
と置く．$\mathcal A$ のスカラー部分と Clifford 積により，
$S$ は $R$ の中心部分環である．Cox 座標での展開は
\begin{equation}
 H^0(Z,L^m)
 =\bigoplus_{j=0}^m c^{m-j}d^j\,k[s,t]_{2m-j}
 \qquad(m\ge0)
 \label{eq:scroll-monomials}
\end{equation}
である．二変数斉次多項式の積が各次数で全射になることから，
$S$ は5つの切断
\[
 cs^2,\quad cst,\quad ct^2,\quad ds,\quad dt
\]
で生成される．一方，式\eqref{eq:explicit-q}の Clifford 関係は
\begin{equation}
 cs^2=z^2,\qquad cst=zw=wz,\qquad ct^2=w^2,
 \qquad ds=\frac{xz+zx}{2},\qquad dt=\frac{xw+wx}{2}
 \label{eq:scalar-generators}
\end{equation}
を与える．従って $S$ 全体が $x,y,z,w$ の生成する部分代数に含まれる．

同様に
\[
 H^0(Z,L^m\otimes\mathcal O_Z(f))
 =\bigoplus_{j=0}^m c^{m-j}d^j\,k[s,t]_{2m+1-j}
\]
なので，次数付き $S$-加群
$\bigoplus_{m\ge0}H^0(Z,L^m\otimes\mathcal O_Z(f))$ は
最初の成分 $H^0(Z,\mathcal O_Z(f))=\langle s,t\rangle$ で生成される．
従って式\eqref{eq:underlying-bundles}の外積分解を用いると，
偶数次数部分は $S$ 上，スカラー $1$ と
\[
 x\wedge y,\quad x\wedge z,\quad x\wedge w,\quad
 y\wedge z,\quad y\wedge w
\]
で生成される．ここに現れる外積は全て次数2の元であり，
例えば $x\wedge y=(xy-yx)/2$ なので
$x,y,z,w$ の積の線形結合である．
奇数次数部分は $S$ 上，$x,y,z,w$ と
\[
 x\wedge y\wedge z,\qquad x\wedge y\wedge w
\]
で生成される．最後の二つは次数3に属し，それぞれ
$x,y,z$ と $x,y,w$ の積を全反対称化した元である．
以上より $R$ の全ての次数が $x,y,z,w$ で生成される．
同時に上の有限生成系は $R$ が $S$ 上有限であることを示す．
$S$ は有限生成可換 $k$-代数なので，$R$ は左右 noetherian である．

最後に一般点で $D_0=\mathcal A\otimes k(Z)$ と置く．
命題\ref{prop:division-degree-three}により $D_0$ は division algebra である．
可逆 bimodule $\mathcal P\otimes k(Z)$ は右 $D_0$-加群として階数1であり，
その tensor 冪の直和は skew polynomial ring と同型になる．
各 $\mathcal P^{\otimes_{\mathcal A}n}$ は中心上のベクトル束なので，
その大域切断は一般点の切断に単射する．従って $R$ はこの
skew polynomial ring の部分代数であり，domain である．
\end{proof}

\subsection{六本の二次関係と三次成分の計算}

命題\ref{prop:generation-all-degrees}より積$R_1\otimes R_1\to R_2$は全射である．
命題\ref{prop:Hilbert-degree-three}から$\dim R_1=4$，$\dim R_2=10$なので，
その核は六次元である．

\begin{lemma}\label{lem:six-quadrics}
上の核は次の独立な六本の二次関係で張られる．
\begin{align}
 r_1&=zw-wz,&
 r_2&=yz+zy,&
 r_3&=yw+wy,\notag\\
 r_4&=2x^2-2z^2-xw-wx,&
 r_5&=xy+yx-2w^2,&
 r_6&=2y^2-2w^2-xz-zx.
 \label{eq:six-quadrics}
\end{align}
$V=R_1$，$W=\langle r_1,\ldots,r_6\rangle$と置くと，二次包絡
$Q=T_k(V)/(W)$は
\[
 \dim_k Q_3=19
\]
を満たす．
\end{lemma}
\begin{proof}
六本の式がClifford関係から従うことは，
式\eqref{eq:explicit-q}と\eqref{eq:scalar-generators}から分かる．
それらの独立性は，順に$zw,yz,yw,x^2,xy,y^2$の係数を見ると従う．
核の次元が六であるからこれで核全体となる．

三次成分は
\[
 Q_3=V^{\otimes3}/(W\otimes V+V\otimes W)
\]
である．$(x_1,x_2,x_3,x_4)=(x,y,z,w)$とし，
三次単項式をこの順の辞書式順序で並べる．
$(r_jx_i)_{j,i}$，$(x_ir_j)_{j,i}$をこの順に並べた
$64\times48$整数行列を$M$とする．
$M$には行列式$512$の$45\times45$小行列が存在し，
また三本の独立な整数右核ベクトルが存在する．
小行列の添字，核ベクトル，および有理数体上の検証コードを
第\ref{sec:reproducibility}節に指定する．
したがって$\operatorname{rank}M\ge45$かつ
$\operatorname{rank}M\le48-3=45$であり，
$\dim Q_3=64-45=19$となる．
\end{proof}

命題\ref{prop:generation-all-degrees}より$Q_3\to R_3$は全射であり，
両辺の次元が十九なので同型である．すなわち，二次関係の三次帰結だけで
反標準環の三次成分の次元に達している．

\subsection{構造対象と偏極を保持した非表示}

標準的なnc $\mathbb P^3$の斉次座標環は，次数一で生成される
二次AS正則代数$G$で$H_G(t)=(1-t)^{-4}$を満たすものとする．
$g\in G_3$が正則正規元なら$B=G/(g)$について
$B_1=G_1$，$B_2=G_2$であり，$G$は二次であるから
\[
 G=T_k(B_1)/\bigl(\ker(B_1\otimes B_1\longrightarrow B_2)\bigr).
\]
従ってその二次包絡の三次成分は必ず$\dim G_3=20$である．
補題\ref{lem:six-quadrics}は，$R$がこのような三次商に次数付き同型でないことを示す．
以下では構造対象を保存しないアーベル圏同値も排除する．
